\chapter{Religions}

\section{Soul as the source Consciousness}

Modern science avoids the word ``soul'' because it cannot be tested, but philosophers, neuroscientists, and even some physicists continue to explore whether consciousness requires something beyond ordinary physical processes. In effect, some modern theories echo older soul-like concepts, even if they do not use the word.

John Eccles argued that consciousness cannot be explained by material processes alone; he proposed a non-physical ``self'' (similar to a soul) interacting with the brain.
Penrose \& Hameroff (Orch-OR theory) suggests that consciousness may arise from quantum processes in neurons and speculated that this might connect with the idea of a soul,
though their theory remains widely disputed.

David Chalmers formulation of the ``hard problem of consciousness'' does not invoke a soul explicitly, but it reopens the possibility by arguing that
subjective experience appears irreducible to brain activity.

Near-death experience researchers (e.g., Pim van Lommel) argue that NDEs indicate consciousness can exist independently of the brain, again resonating with
traditional soul concepts.

Thus, while the term ``soul'' has largely dropped out of scientific vocabulary, the debate surrounding consciousness often returns to closely related conceptual territory.

\subsection*{NDEs as Evidence for Soul}

While every story appears somewhat unique, researchers like Pim van Lommel have identified several common core elements that people report during these episodes:

\begin{itemize}
\item \textbf{Out-of-Body Experiences (OBEs):} The sensation of floating above one’s own body and observing medical procedures or surroundings.
\item \textbf{The "Tunnel":} A feeling of being drawn through a dark space toward a brilliant, warm light.
\item \textbf{Life Review:} A panoramic replay of one’s life events, often felt from the perspective of others involved.
\item \textbf{Intense Peace:} An overwhelming sense of love, euphoria, and the absence of pain.
\item \textbf{The "Border":} Reaching a point of no return where they are told (or decide) to go back to their body.
\end{itemize}

But how reliable are these findings?

Pim van Lommel is a Dutch cardiologist who published a landmark study in The Lancet (2001). His work is significant because instead of just
listening to old stories, he studied patients who were resuscitated in hospitals in real-time. He found that some patients reported clear,
structured consciousness at a time when their brains had no measurable electrical activity (clinical death).
Van Lommel argues that the brain might act more like a receiver for consciousness rather than the producer of it. If the radio (the brain) is broken,
the broadcast (consciousness) still exists; one just can’t hear it. Unlike hallucinations, NDEs almost always permanently change a person’s personality
and remove their fear of death.

Science admits that consciousness is far more resilient than once thought. The lights don't go out instantly. We know the house stays "powered" for quite a
while after the main generator fails. Whether that power comes from a backup battery (the brain) or a different power grid entirely (the soul) is
still the ultimate mystery.


\subsection*{New Testament and Christianity}

Should we turn to God? Should we replace our physics books with the Bible?

According to Jesus (Matthew 10:28):
``Do not fear those who kill the body but cannot kill the soul.''

The New Testament indirectly links many aspects of consciousness
to the concept of the soul. Our conscious decisions determine
what happens to our soul once we die. The soul is described as the
immaterial and eternal part of a human being that is distinct from
the physical body.

Indeed, soul, just like consciousness and pain, appear to live in a domain
that is not physical. It is something that cannot be touched, weighted,
or directly measured by any measuring device.

How reliable is the New Testament? How strong of a case do
science and archaeological findings make to support the stories
within it? Is there any evidence to support that a person named
Jesus ever lived?

It turns out there is no single piece of direct archaeological
evidence for Jesus whatsoever. We just have to believe the story
put forward in the Gospels.

The Gospels are, however, based on a large number of ancient
archaeological documents written in Greek. Thousands of manuscripts
have been found, and new ones are discovered every year.

None of these manuscripts are original, but merely copies of
copies. Due to the manual copying methods used back then (copy
machines were invented much later), none of the found manuscripts
is identical. They all contain errors and differences.
However, by comparing the numerous copies found, researchers have
managed to restore the original document.

From the four Gospels, three would seem to tell essentially the
same story. These are Matthew's, Mark's, and Luke's Gospels, and
they are known as the \texttt{Synoptic} Gospels. There are 661
verses in Mark's Gospel, from which 607 are also included in
Matthew's Gospel, and 360 are included in Luke's Gospel. Matthew
and Luke have 230 common verses which, however, are not included
in Mark. Because of this, it would seem that Matthew and Luke are
based on Mark, as well as on some yet unknown source that is
called the ``Q'' document. The name ``Q'' comes from the German
word ``Quelle,'' which means ``source.'' Q is thought to be a
collection of sayings and teachings of Jesus that were shared by
Matthew and Luke but not found in Mark.

The Gospel of Mark is believed to have been written between
60--70 AD. Science has managed to pinpoint the dating of the
manuscripts with astonishing accuracy by considering many
factors, for example handwriting style, paleographic analysis,
and historical references.

So the fact is that the New Testament draws its foundation from
thousands of ancient documents. Rejecting their authenticity is
akin to disputing the existence of dinosaurs despite the
continual discovery of new fossils each year. Moreover, early
dating provides the New Testament with a degree of credibility,
offering testimony that, while not necessarily from eyewitnesses,
still carries weight.

\subsection*{Doubt}

However, a couple of concerns arise in the mind of an average
programmer.

All the manuscripts seem to be based on just two original
sources: the Gospel document and the yet-to-be-found Q-document
(as per the ``Two-Source Hypothesis''). How can one be certain
that the original text was not written by someone suffering from
a wild imagination, at the very least? How can one determine that
the authors of the original texts did not exaggerate to some
degree? How can we discern whether these authors were simply
storytellers of their time?

Furthermore, Jews do not seem to believe in Jesus as the Messiah.
According to Jewish tradition several reasons exist. Jesus did
not fulfill the messianic prophecies nor embody the personal
qualifications of the Messiah. He did not build the Third Temple
nor gather all Jews back to the Land of Israel. He also failed to
spread knowledge of the God of Israel so that humanity would be
united as one. The God of Israel is definitely not king over all
the world.

Jesus' teachings and the doctrines associated with Christianity,
including the concept of the Trinity and the divinity of Jesus,
are also in serious contradiction with Jewish theological
beliefs. Judaism emphasizes monotheism and the unity of God,
rejecting the notion of Jesus as a divine figure.

What worries an average programmer is that the God worshipped by
Christians and Jews is the same God. Both religions trace their
roots back to the Hebrew Bible (Old Testament) and acknowledge
God as the one and only.

The fact that Jews do not believe in Jesus, the son of God, as
the Messiah therefore feels like a serious matter. It is the very
soul of an average programmer that is at stake. Jesus was a Jew,
and if Jews themselves do not believe in him, then why should an
average programmer do so?

Are we Christians certain that we are on the right path? Are we
confident that we are headed toward heaven instead of eternal
suffering in the fires of hell?

And what about other religions, many of which seem even more
incompatible with Christianity? Do they fall into the category
of false religions, with believers whose unfortunate fate one
can only regret?

